\clearpage
\appendix
\section{Appendix}
\subsection{Implement Details}

\noindent\textbf{B-spline action representation. }
We represent action trajectories using cubic B-splines with degree $k=3$.
In simulation, we follow standard benchmark settings and use a delta end-effector action space.
In real-world experiments, we use absolute joint-position actions.

\noindent\textbf{Adaptive b-spline fitting from demonstrations. }
We fit a B-spline to each demonstration trajectory using the adaptive knot insertion procedure in Alg.~\ref{alg:spline-compress}.
For most tasks, we use a small fitting tolerance $\varepsilon=0.002$ to preserve high-precision behaviors.
For Push-T, we use $\varepsilon=1$ because the action magnitude is substantially larger than in the other benchmarks, with values ranging roughly from $0$ to $512$.

\noindent\textbf{Model training of \our.} 
We train \our model to predict a fixed-size parameter segment consisting of 16 knots and their corresponding control points, including six boundary-support knots.
Because observations are not aligned to knot or control-point indices, at each observation step, the model predicts the nearest future segment containing the next 16 knots and control points.
In simulation, all models use the RoboMimic visual encoder.
In real-world experiments, diffusion uses the same visual encoder, whereas the regression uses a DINOv2 visual encoder.

\noindent\textbf{Model training of action-chunk policies.} 
For a fair comparison, all action-chunk baseline policies (Diffusion/Regression) are also trained to predict 16 future actions.
Each baseline is trained for the same number of epochs as its B-spline counterpart.
The visual encoder is the same as \our.

\noindent\textbf{DemoSpeedup.} We implement the DemoSpeedup based on the diffusion policy. 
The method uses two key hyperparameters, which are the downsampling rates for the precision and non-precision phases.
We found that downsampling the precision phase by a factor $\ge 2$ prevents the policy from completing the task reliably.
We therefore set the precision-phase downsampling rate to 1.
For the non-precision phase, we use a downsampling rate of 4 for \pickcube\ and \cleantable.
We use 2 for \stackunstack\ because it requires higher accuracy.

\noindent\textbf{Controller.}
In simulation, actions are delta end-effector commands.
We sample B-spline actions at 100\,Hz, but we do not apply segment alignment because it does not have a clear physical interpretation for delta actions.
In real-world experiments, actions are absolute joint positions.
We sample actions from each predicted B-spline segment at 100\,Hz.
For phase alignment, we set the search window to $[0,\lambda T^{\mathrm{inf}}]$ with $\lambda T^{\mathrm{inf}} = u_{\max}(S)/2$, where $u_{\max}(S)$ denotes the segment horizon in the spline parameter domain.

% \noindent\textbf{PD controller with feedforward terms.}
% In Eq.~\ref{eq:ff} of the main text, $\boldsymbol{q}, \dot{\boldsymbol{q}}$ are the measured joint positions and velocities; $\boldsymbol{q}_d, \dot{\boldsymbol{q}}_d, \ddot{\boldsymbol{q}}_d$ are the desired joint positions, velocities, and accelerations, computed by evaluating the B-spline trajectory and its first and second derivatives at the current sample time. $\boldsymbol{K}_p, \boldsymbol{K}_d$ are diagonal positive-definite gain matrices; $\boldsymbol{M}(\boldsymbol{q})$ is the manipulator inertia matrix; $\boldsymbol{C}(\boldsymbol{q}, \dot{\boldsymbol{q}})$ captures Coriolis and centripetal effects; and $\boldsymbol{g}(\boldsymbol{q})$ denotes gravitational torques.

% \subsection{Discussion}
% % not sure if we need to add this

% \noindent\textbf{Relationship to DemoSpeedup.}
% DemoSpeedup partitions trajectories into precision and non-precision phases using action entropy and applies a higher downsampling rate to the non-precision phase.
% It requires training a separate policy for each choice of downsampling rates, which is expensive when sweeping speed settings.
% In contrast, \our enables speed changes at inference time via temporal scaling of a continuous trajectory.
% Besides, the two approaches are actually complementary.
% Entropy-based phase-dependent speed scheduling can also be applied on top of \our.

% \noindent\textbf{Relationship to SAIL.} SAIL~\cite{arachchige2025sail} also focuses on enabling faster-than-demonstration execution of visuomotor policies. It proposed a full-stack system, which is orthogonal to our learning representation. The main focus of our work is to propose a new action representation, which is more friendly for the robot to finish tasks faster. 

% \noindent\textbf{Relationship to post-hoc spline smoothing.}
% A natural simplified variant is to fit a B-spline online to each waypoint chunk predicted by a standard policy and then execute it with the same high-rate controller pipeline, and segment align scheme as \our.
% While this might speed up simple tasks, it remains limited by the waypoint representation: the policy still learns discrete samples at a fixed temporal resolution.
% In contrast, \our directly predicts a \emph{non-uniform} spline parameterization, allocating more knots to rapidly changing phases and fewer to smooth segments.
% Moreover, chunk-wise online fitting is inherently local, whereas fitting splines from full demonstrations can capture global structure during preprocessing, with more room for optimization.
% Empirically, \our improves success rates at the original speed on most tasks compared to the discrete action representation. We hypothesize that non-uniform knot placement contributes to this effect.

\noindent\textbf{Knot validity projection.}
A valid B-spline requires a nondecreasing knot vector, which standard imitation-learning regressors do not enforce. In practice predicted knots are almost always monotonic; for any knot $u_i$ that violates the constraint we apply $u_i \leftarrow \max(u_i,\, u_{i-1} + \delta)$ with a small constant $\delta > 0$, restoring validity with minimal perturbation.

\noindent\textbf{Inference techniques.} During inference, pipelined execution with segmentment alignment is applied. Here we detailed the complete procedure in Algorithm~\ref{alg:pipeline-inference}.

\begin{table}[h]
\centering
\setlength{\tabcolsep}{10.0pt}
\begin{tabular}{c|cc}
\toprule
Push-T & Score & Compression Ratio \\
\midrule
$\varepsilon{=}0.5$ & 0.52 & 1.12 \\
$\varepsilon{=}1$   & 0.65 & 1.34 \\
$\varepsilon{=}2$   & 0.62 & 1.73 \\
$\varepsilon{=}4$   & 0.63 & 2.36 \\
$\varepsilon{=}8$   & 0.65 & 3.34 \\
\bottomrule
\end{tabular}
\caption{\textbf{Ablation on fitting error tolerance.}
We report the success rate and the compression ratio under different maximum fitting errors $\varepsilon$.}
\label{tab:pushT_fit_error}
\end{table}


\begin{algorithm}[h]
% \setlength{\intextsep}{0pt}
% \setlength{\columnsep}{6pt}
% \begin{wrapfloat}{algorithm}[24]{r}{0.5\textwidth}
\caption{Inference pipeline}
\label{alg:pipeline-inference}
\textbf{Input:} policy $\pi$, speedup factor $m$, inference budget $\delta^{\text{budget}}$ \\
\textbf{State:} current B-spline segment $S$, segment start time $t_0$, in-flight inference flag \quad in\_flight
\begin{enumerate}[leftmargin=1.2em]
\item \textbf{Initialize:} $S \leftarrow \pi(o_{\text{recent}})$,\quad $t_0 \leftarrow \text{now}()$,\quad in\_flight $\leftarrow$ false.
\item \textbf{Loop (control at high rate):}
\begin{enumerate}[leftmargin=1.2em]
\item $t_{\text{real}} \leftarrow \text{now}() - t_0$, \quad $u \leftarrow m \cdot t_{\text{real}}$ \hfill(\textit{Speedup})
\item Execute $a \leftarrow S(u)$ \hfill(\textit{Sample})
\item \textbf{If} $(u_{\max}(S) - u) < \delta^{\text{budget}}$ and in\_flight = false:
\begin{enumerate}[leftmargin=1.2em]
\item Launch asynchronous inference to obtain a new segment $S_{\text{new}} \leftarrow \pi(o_{\text{recent}})$
\item Set in\_flight $\leftarrow$ true
\end{enumerate}
\item \textbf{On inference completion:} let $T^{\text{inf}}$ be the measured inference latency.
\begin{enumerate}[leftmargin=1.2em]
\item Compute alignment time $t^\star$ by Eq.~\eqref{eq:align-time} (using $S_{\text{new}}$ and $a_{\text{last}}$)
\item Swap $S \leftarrow S_{\text{new}}$
\item Update $t_0 \leftarrow \text{now}() - \frac{t^\star}{m}$ \hfill(\textit{Segment align})
\item Set in\_flight $\leftarrow$ false
\end{enumerate}
\end{enumerate}
\end{enumerate}
% \end{wrapfloat}
\end{algorithm}

\subsection{Additional experiments results}


\noindent\textbf{Ablation on fitting error tolerance.}
Representing actions with B-splines introduces a fitting tolerance $\varepsilon$.
This typically yields fewer knot and control point pairs than the original demonstrations.
Tab.~\ref{tab:pushT_fit_error} reports success rate and compression ratio under different tolerances.
We define the compression ratio as the number of original actions divided by the number of knot--control-point pairs.
As $\varepsilon$ increases, the compression ratio increases substantially while success rate remains relatively stable.
This suggests that performance is not highly sensitive to $\varepsilon$ within a reasonable range, and a moderately larger tolerance can reduce computational cost.

% \subsection{Details of B-Spline Definition and Properties}
% \subsection{Physical meaning of B-Spline}


% \begin{equation*}
% \begin{aligned}
% B_{i,3}(x) &=
% \frac{x - t_i}{t_{i+3} - t_i} B_{i,2}(x) +
% \frac{t_{i+4} - x}{t_{i+4} - t_{i+1}} B_{i+1,2}(x) \\
% B_{i,2}(x) &=
% \frac{x - t_i}{t_{i+2} - t_i} B_{i,1}(x) +
% \frac{t_{i+3} - x}{t_{i+3} - t_{i+1}} B_{i+1,1}(x) \\
% B_{i,1}(x) &=
% \frac{x - t_i}{t_{i+1} - t_i} B_{i,0}(x) +
% \frac{t_{i+2} - x}{t_{i+2} - t_{i+1}} B_{i+1,0}(x) \\
% &=\frac{x - t_i}{t_{i+1} - t_i}
% \end{aligned}
% \end{equation*}

% \[
% h_r := t_{i+r+1}-t_{i+r},\qquad r=0,1,2,3.
% \]

% \[
% B_{i,3}(x)=
% \begin{cases}
% \displaystyle \frac{(x-t_i)^3}{h_0\,(h_0+h_1)\,(h_0+h_1+h_2)},
% & t_i\le x<t_{i+1},\\[12pt]

% \displaystyle \frac{A_3 u^3 + A_2 u^2 + A_1 u + A_0}{D_1},
% & t_{i+1}\le x<t_{i+2},\quad (u=x-t_{i+1}),\\[12pt]

% \displaystyle \frac{C_3 v^3 + C_2 v^2 + C_1 v + C_0}{D_2},
% & t_{i+2}\le x<t_{i+3},\quad (v=x-t_{i+2}),\\[12pt]

% \displaystyle \frac{(t_{i+4}-x)^3}{h_3\,(h_2+h_3)\,(h_1+h_2+h_3)},
% & t_{i+3}\le x<t_{i+4},\\[10pt]
% 0,& \text{otherwise}.
% \end{cases}
% \]

% % --- Piece on [t_{i+1}, t_{i+2}) ---
% \[
% D_1 = h_1\,(h_0+h_1)\,(h_1+h_2)\,(h_0+h_1+h_2)\,(h_1+h_2+h_3),
% \]
% \[
% S := h_1^2 + 2h_1h_2 + h_1h_3 + h_2^2 + h_2h_3,
% \]
% \[
% \begin{aligned}
% A_3 &= -\Big(h_0^2 + 3h_0h_1 + 2h_0h_2 + h_0h_3 + 3h_1^2 + 4h_1h_2 + 2h_1h_3 + h_2^2 + h_2h_3\Big),\\
% A_2 &= 3h_1\,S,\\
% A_1 &= 3h_0h_1\,S,\\
% A_0 &= h_0^2h_1\,S.
% \end{aligned}
% \]

% % --- Piece on [t_{i+2}, t_{i+3}) ---
% \[
% D_2 = h_2\,(h_1+h_2)\,(h_2+h_3)\,(h_0+h_1+h_2)\,(h_1+h_2+h_3),
% \]
% \[
% \begin{aligned}
% C_3 &= h_0h_1 + 2h_0h_2 + h_0h_3 + h_1^2 + 4h_1h_2 + 2h_1h_3 + 3h_2^2 + 3h_2h_3 + h_3^2,\\
% C_2 &= -3\Big(h_0h_2^2 + h_0h_2h_3 + 2h_1h_2^2 + 2h_1h_2h_3 + 2h_2^3 + 3h_2^2h_3 + h_2h_3^2\Big),\\
% C_1 &= -3\Big(h_0h_1h_2^2 + h_0h_1h_2h_3 + h_1^2h_2^2 + h_1^2h_2h_3\Big)
%       +3\Big(h_2^4 + 2h_2^3h_3 + h_2^2h_3^2\Big),\\
% C_0 &= \;h_0h_2^4 + 2h_1h_2^4 + 2h_1^2h_2^3 + 2h_0h_1h_2^3 + 2h_0h_2^3h_3 + 4h_1h_2^3h_3\\
% &\quad + 3h_0h_1h_2^2h_3 + 3h_1^2h_2^2h_3 + h_0h_2^2h_3^2 + 2h_1h_2^2h_3^2 + h_0h_1h_2h_3^2 + h_1^2h_2h_3^2.
% \end{aligned}
% \]

% \[
% a(x) = c_{i-3} B_{i-3}(x) + c_{i-2} B_{i-2}(x)  + c_{i-1} B_{i-1}(x) + c_iB_i(x) (t_i \le x \le t_{i+1} )
% \]

% Let $(t_i, $

% \begin{equation}
% \tau = K_p(q_d - q) + K_d(\dot q_d - \dot q) .
% \end{equation}

% \begin{equation}
% M(q)\ddot q + C(q,\dot q)\dot q + g(q) = \tau .
% \end{equation}

% \begin{equation}
% \tau_{\mathrm{ff}} =
% M(q_d)\ddot q_d + C(q_d,\dot q_d)\dot q_d + g(q_d) .
% \end{equation}

% \begin{equation}
% \tau_{\mathrm{apply}} = \tau_{\mathrm{ff}} + K_p(q_d - q) + K_d(\dot q_d - \dot q) .
% \end{equation}

% Let $I=[a,b]$ and let $S^p_{\tau}$ denote the spline space of degree $p$
% over a knot partition $\tau$ on $I$ (equivalently representable using a B-spline basis).
% A classical result in spline approximation states that, as the mesh size $|\tau|\to 0$,
% these spline spaces are dense in $C(I)$ in the uniform norm.
% Consequently, for any continuous parametric curve $\gamma:I\to\mathbb{R}^d$
% and any $\varepsilon>0$, there exists a B-spline curve $c:I\to\mathbb{R}^d$ such that
% \[
% \|\gamma-c\|_{\infty}
% := \sup_{t\in I}\|\gamma(t)-c(t)\|_2
% < \varepsilon.
% \]

% For the zeroth degree,
% \begin{equation}
% B_{i,0}(u) =
% \begin{cases}
% 1, & t_i \le u < t_{i+1},\\
% 0, & \text{otherwise},
% \end{cases}
% \end{equation}
% and for $k \ge 1$,
% \begin{equation}
% B_{i,k}(u) = \alpha_{i,k}(u)\, B_{i,k-1}(u) \;+\; \beta_{i,k}(u)\, B_{i+1,k-1}(u),
% \end{equation}
% where the blending weights are
% \begin{equation}
% \alpha_{i,k}(u) = \frac{u - t_i}{t_{i+k} - t_i}, 
% \qquad
% \beta_{i,k}(u) = \frac{t_{i+k+1} - u}{t_{i+k+1} - t_{i+1}}.
% \end{equation}
% By convention, if a denominator is zero (which can occur with repeated knots), the corresponding
% fraction is taken to be zero.

% \section{}

% \begin{figure}
%     \begin{minipage}[t]{0.5\linewidth}
%     \vspace{20pt}
%     \centering
%     % \includegraphics[width=0.98\linewidth]{images/data-fitting-3.pdf}
%     \includegraphics[width=1.\linewidth]{images/data-fitting-4.pdf}
%     \end{minipage}
%     \begin{minipage}[t]{0.48\linewidth}
%     \vspace{5pt}
%     % \vspace{-90pt}
%     \caption{\textbf{Fitting a B-spline.} We show a concept visualization of how we fit a B-spline. Given a sequence of waypoints, we apply Algorithm~\ref{alg:spline-compress} to fit a B-spline, which adaptively selects knot locations and estimates control points such that the fitting error remains below a prescribed threshold.}
%     \label{fig:data-fitting}
%     \end{minipage}
%     \vspace{-20pt}
% \end{figure}


% \begin{algorithm}[t]
%     \caption{Adaptive B-spline Trajectory Fitting}
%     \label{alg:spline-compress}
    
%     \textbf{Input:}
%     Trajectory samples $\{(t_i, x_i)\}_{i=1}^N$, 
%     error tolerance $\varepsilon$,
%     (optional) maximum number of knots $K_{\max}$\\
%     \textbf{Output:}
%     Knot set $\mathcal{K}$ and spline $g$
    
%     \begin{enumerate}[leftmargin=1.2em]
    
%     \item \textbf{Initialize:}
%           Set $\mathcal{K} \gets \{t_1, t_N\}$ (boundary knots only).
    
%     \item \textbf{Repeat until convergence or knot budget exhausted:}
    
%         \begin{enumerate}[leftmargin=1.2em]
    
%         \item \textbf{spline fitting.}  
%               Fit a spline $g(\cdot)$ on $(t_i, x_i)$ using current knots $\mathcal{K}$.
    
%         \item \textbf{Error evaluation.}  
%               Compute reconstruction errors
%               \[
%                   e_i = |g(t_i) - x_i|, 
%                   \quad E = \max_i e_i .
%               \]
%               If $E \le \varepsilon$, stop.
    
%         \item \textbf{Interval error accumulation.}  
%               Partition the domain into intervals $[k_j, k_{j+1})$ for all consecutive knots $k_j \in \mathcal{K}$.
%               For each interval compute
%               \[
%                   E_j = \sum_{i : t_i \in [k_j, k_{j+1})} e_i^2 .
%               \]
    
%         \item \textbf{Select the worst interval.}
%               \[
%                   j^\star = \arg\max_j E_j .
%               \]
    
%         \item \textbf{Insert a new knot.}  
%               Choose the sample point with the largest error in this interval:
%               \[
%                   t_{\text{new}}
%                   = \arg\max_{t_i \in [k_{j^\star},\, k_{j^\star+1})} e_i .
%               \]
%               Update $\mathcal{K} \gets \mathcal{K} \cup \{t_{\text{new}}\}$.
    
%         \end{enumerate}
    
%     \end{enumerate}
    
%     \textbf{return} $\mathcal{K}$ and spline $g$
    
%     \end{algorithm}

% \section{More experiments}